\section{Structured Approach to System Architecture and Module Topology}

\subsection{Structured Approach to Express Application Topology}

The Express app composes JSON body parsing, targeted CORS handling, API route registration, admin reset tooling, static upload serving, and SPA-like fallback routing for public assets \cite{repo-server-app,express-docs}. The API namespace is normalized through \texttt{API\_PREFIX = /api} inside shared HTTP helpers \cite{repo-security}.

\begin{figure}[h]
\centering
\begin{tikzpicture}[
  node distance=8mm and 10mm,
  box/.style={draw, rounded corners, align=center, minimum width=3.2cm, minimum height=0.9cm, font=\small},
  data/.style={draw, cylinder, shape border rotate=90, aspect=0.25, align=center, minimum height=1.0cm, minimum width=2.2cm, font=\small},
  arrow/.style={-Latex, thick}
]
\node[box] (client) {Browser Client\\(public pages / React client)};
\node[box, below=of client] (app) {Express App\\\texttt{server/app.js}};
\node[box, below left=of app, xshift=-10mm] (modules) {Domain Modules\\auth, classes, assignments,\\messages, discussions, events,\\resources, uploads, AI, admin};
\node[data, below=of modules] (jsondb) {\texttt{learning-os.json}};
\node[box, below right=of app, xshift=8mm] (cloudinary) {Cloudinary Upload\\Gateway};
\node[box, right=of cloudinary, xshift=5mm] (firebase) {Firebase Migration\\Script Path};

\draw[arrow] (client) -- node[right,font=\scriptsize]{HTTP / JSON / multipart} (app);
\draw[arrow] (app) -- (modules);
\draw[arrow] (modules) -- (jsondb);
\draw[arrow] (modules.east) -- ++(1.3,0) |- (cloudinary.west);
\draw[arrow] (jsondb.east) -- ++(1.3,0) |- (firebase.west);
\end{tikzpicture}
\caption{Repository-implemented architecture topology.}
\label{fig:architecture-topology}
\end{figure}

\subsection{Structured Approach to Backend Domain Boundaries}

\texttt{server/modules/index.js} registers 12 module groups in a deterministic order: authentication, admin, dashboard, classes, assignments, discussions, messages, notifications, events, resources, AI helpers, and upload handling \cite{repo-server-modules}. This achieves bounded routing concerns without introducing cross-process microservices.

\begin{table}[h]
\centering
\caption{Module Boundary and Core Responsibility Mapping}
\begin{tabularx}{\linewidth}{>{\raggedright\arraybackslash}p{3.2cm} >{\raggedright\arraybackslash}X >{\raggedright\arraybackslash}p{3.1cm}}
\toprule
Module & Core Responsibility & Access Pattern \\
\midrule
Auth & Health/meta endpoints, credential validation, session token issuance & Public + authenticated \\
Admin & User lifecycle and role management & Admin only \\
Dashboard & Snapshot access, widget preference updates, analytics overview & Authenticated \\
Classes & Class create/join/archive lifecycle & Role-gated write operations \\
Assignments & Assignment retrieval and creation & Teacher/admin create \\
Upload & Resource/submission file handling and grading updates & Mixed role controls \\
Notifications & Inbox fetch, mark-read mutation, SSE unread stream & Authenticated \\
AI & Summaries, quiz generation, feedback synthesis & Input-validated POST endpoints \\
\bottomrule
\end{tabularx}
\end{table}

\subsection{Structured Approach to Architecture Reality Check}

Architecture documentation files describe a planned stack using PostgreSQL, Redis, and S3-compatible storage. Current executable code persists state in JSON files and uses Cloudinary for file object storage. Therefore, the implemented architecture is best classified as a local-file-backed modular monolith with optional cloud upload and optional Firebase migration utilities, rather than a distributed data platform \cite{repo-store,repo-upload,firebase-admin-docs}.
